\section{Proposed Method}
\label{sec:method}

We study \emph{asymmetric} self-supervised objectives for latent-space JEPA:
instance-level contrastive learning on a masked predictive path, paired with
information-maximization regularizers on an augmentation path.
The goal is to retain the strong local geometry of InfoNCE while improving
global structure for linear classification.

\subsection{Architecture and views}
\label{sec:architecture}

We use a single online encoder $f_\theta$ and a lightweight predictor
$g_\phi$ (no EMA teacher).
Given an image $x$, we construct three views:
\begin{itemize}
  \item \textbf{Clean} $\tilde{x}$: standard random crop and horizontal flip.
  \item \textbf{Corrupt} $x_c$: multi-block masking ($60\%$ of $4{\times}4$ patches
        zeroed, I-JEPA style~\cite{assran2023ijepa}).
  \item \textbf{Augmented} $x_a$: color jitter (brightness, contrast, saturation, hue).
\end{itemize}
All views are encoded into $d$-dimensional latents ($d{=}256$ in our experiments):
\[
  z = f_\theta(\tilde{x}), \quad
  z_c = f_\theta(x_c), \quad
  z_a = f_\theta(x_a).
\]
The predictor maps the corrupt latent to the clean target space:
$\hat{z} = g_\phi(z_c)$.
At evaluation, only $f_\theta(\tilde{x})$ is used; $g_\phi$ is discarded.

\subsection{Dual InfoNCE baseline}
\label{sec:baseline}

Our starting point replaces cosine JEPA regression with batch-wise InfoNCE
(SimCLR-style~\cite{chen2020simclr}) on two alignment pairs:
\begin{align}
  \mathcal{L}_{\mathrm{jepa}}
  &= \mathrm{InfoNCE}(\hat{z}, \sg(z); \tau), \\
  \mathcal{L}_{\mathrm{aug}}
  &= \mathrm{InfoNCE}(z_a, \sg(z); \tau),
\end{align}
where $\sg(\cdot)$ stops gradients on the clean encoder branch and
$\tau{=}0.1$ is the temperature.
The total loss is
$\mathcal{L}_{\mathrm{dual}} =
 w_j \mathcal{L}_{\mathrm{jepa}} + w_a \mathcal{L}_{\mathrm{aug}}$
with $w_j = w_a = 1$.
This setup yields strong k-NN accuracy but weak linear-probe performance,
consistent with instance-discrimination objectives that emphasize local
neighborhood structure over globally linearly separable representations.

For one image $i$ in a batch of size $B$, InfoNCE on queries $q_i$ and keys
$k_j = \sg(z_j)$ is
\begin{equation}
  \mathrm{InfoNCE}(q, k; \tau)
  = -\log \frac{\exp(\cos(q_i, k_i)/\tau)}
  {\sum_{j=1}^{B} \exp(\cos(q_i, k_j)/\tau)},
\end{equation}
with $\cos(\cdot,\cdot)$ the cosine similarity of $\ell_2$-normalized vectors.

\subsection{Asymmetric hybrids}
\label{sec:hybrid}

We keep $\mathcal{L}_{\mathrm{jepa}}$ on the mask path and \emph{replace}
$\mathcal{L}_{\mathrm{aug}}$ with explicit spread / invariance terms on
$(z_a, z)$, yielding two variants that share the same structure but differ in
the second regularizer.

\paragraph{Shared first term.}
Both hybrids optimize
\begin{equation}
  \mathcal{L}_{\mathrm{mask}}
  = w_j \,\mathrm{InfoNCE}(\hat{z}, \sg(z); \tau),
  \qquad w_j = 1.
\end{equation}

\paragraph{Hybrid A: MSE + variance + covariance.}
On the augmentation path we apply three decomposed terms (same building blocks
as VICReg~\cite{bardes2022vicreg}, but applied asymmetrically and \emph{not}
as a standalone VICReg objective):
\begin{align}
  \mathcal{L}_{\mathrm{inv}} &= \| z_a - z \|_2^2, \\
  \mathcal{L}_{\mathrm{var}} &= \sum_{v \in \{a,\,\mathrm{clean}\}}
    \mathrm{VarReg}(z_v), \\
  \mathcal{L}_{\mathrm{cov}} &= \sum_{v \in \{a,\,\mathrm{clean}\}}
    \mathrm{CovReg}(z_v),
\end{align}
where $\mathrm{VarReg}$ is a hinge on per-dimension batch standard deviation
(target $\ge 1$) and $\mathrm{CovReg}$ penalizes off-diagonal covariance
entries.
The full objective is
\begin{equation}
  \mathcal{L}_{\mathrm{mvc}}
  = \mathcal{L}_{\mathrm{mask}}
  + \lambda_{\mathrm{inv}} \mathcal{L}_{\mathrm{inv}}
  + \lambda_{\mathrm{var}} \mathcal{L}_{\mathrm{var}}
  + \lambda_{\mathrm{cov}} \mathcal{L}_{\mathrm{cov}}.
\end{equation}
We report two weightings: \textbf{unit scale}
($\lambda_{\mathrm{inv}} = \lambda_{\mathrm{var}} = \lambda_{\mathrm{cov}} = 1$)
and \textbf{VICReg scale}
($25/25/25$, matching prior VICReg runs in our codebase).

\paragraph{Hybrid B: MSE + SIGReg.}
Alternatively, we replace $\mathcal{L}_{\mathrm{var}}$ and
$\mathcal{L}_{\mathrm{cov}}$ with Sketched Isotropic Gaussian Regularization
(SIGReg) from LeJEPA~\cite{balestriero2025lejepa}:
\begin{equation}
  \mathcal{L}_{\mathrm{sig}}
  = \mathcal{L}_{\mathrm{mask}}
  + \lambda_{\mathrm{inv}} \mathcal{L}_{\mathrm{inv}}
  + \lambda_{\mathrm{sig}} \,\mathrm{SIGReg}\big([z_a; z]\big),
\end{equation}
where $[z_a; z]$ concatenates both branches along the batch dimension and
SIGReg matches random 1D projections to a standard Gaussian via an
Epps--Pulley statistic.
We use $\lambda_{\mathrm{inv}} = \lambda_{\mathrm{sig}} = 1$ unless noted.

\paragraph{Design rationale.}
Dual InfoNCE applies instance discrimination twice (mask and augmentation),
which encourages tight local neighborhoods but can leave class structure
poorly organized for linear probes.
Our asymmetric design assigns \emph{contrastive} learning to the JEPA path
(where masking creates a non-trivial predictive task) and \emph{spread /
invariance} to the augmentation path (where the task is view matching without
batch negatives).
This follows the spirit of JEPA + redundancy-reduction hybrids such as
C-JEPA~\cite{mo2024cjepa}, but differs in (i) using InfoNCE rather than
patch-level $\ell_2$ regression on the mask path, and (ii) applying spread
regularizers on $(z_a, z)$ rather than on corrupt/clean encoder pairs alone.

\subsection{Implementation summary}
\label{sec:implementation}

\begin{itemize}
  \item \textbf{Backbone:} SmallResNet encoder ($\sim$2.8M parameters including
        predictor) on CIFAR-100 ($32{\times}32$).
  \item \textbf{Optimization:} AdamW, lr $= 3{\times}10^{-4}$, cosine schedule,
        batch size $256$, $100$ epochs, seed $42$.
  \item \textbf{Evaluation:} frozen encoder; k-NN@20 (cosine, normalized
        features) and linear probe (SGD, 50 epochs on train features).
\end{itemize}
